针对圆柱药材预热、恒温干燥及失水收缩过程，建立一维轴对称热质耦合模型。四问均在表面热平衡中计入实际水分通量乘汽化潜热，并区分药材干基含水率与空气湿度比，以内部供水能力和气膜传质能力共同限制蒸发。物性严格分别采用题给附录2、3、4；通过材料坐标变换处理收缩并保持干物质量守恒，以有限体积法和BDF隐式积分求解。

问题一使用附件实测环境，1800~s时中心、表面温度分别为\QOneCenterT{}、\QOneSurfaceT{}~$^\circ$C，含水率分别为\QOneCenterC{}、\QOneSurfaceC{}~kg/kg。问题二建立全程变物性模型，前4~h保留附件输入，后续主方案依据文献选择\EnvTempC{}~$^\circ$C、\EnvRHPercent{}\%相对湿度；3~h时中心温度为\QTwoCenterT{}~$^\circ$C、含水率为\QTwoCenterC{}~kg/kg。

问题三以全域最大含水率低于0.15~kg/kg为判据，得到连续达标时长\QThreeEndH{}~h。问题四结合实测半径及附录4物性，得到\QFourEndH{}~h，终点半径\QFourEndRadiusCm{}~cm。另对附件升温曲线拟合渐近温度，构造\FitTempC{}~$^\circ$C、20\%相对湿度的第二方案，其问题三、四时长分别为\FitQThreeEndH{}、\FitQFourEndH{}~h；该温度依赖外推模型，不能视为已观测到的稳态。

对新模型开展空间加密、时间积分交叉、守恒收支及工况对照，并生成四个完整结果工作簿。药材种类与解吸数据缺失导致表面供水闭合仍有结构不确定性，数值收敛不等于物种实验验证。

\keywords{药材干燥；蒸发潜热；热质耦合；移动边界；有限体积法}
